\documentclass[11pt,a4paper]{article}
\usepackage{fontspec,amsmath,amssymb,amsthm,microtype}
\setmainfont{lmroman10-regular.otf}[BoldFont=lmroman10-bold.otf,ItalicFont=lmroman10-italic.otf,BoldItalicFont=lmroman10-bolditalic.otf]
\setmonofont{lmmono10-regular.otf}[Scale=MatchLowercase]
\usepackage[margin=25mm]{geometry}
\usepackage[unicode,hidelinks]{hyperref}
\setlength{\emergencystretch}{3em}
\setlength{\parskip}{0.4em}
\setlength{\parindent}{0pt}
\allowdisplaybreaks[2]
\raggedbottom
\title{\texorpdfstring{Sensitivity to the proximal parameter}{Sensitivity to the proximal parameter}}
\author{LangevinLean mathematical exposition}
\date{}
\begin{document}\maketitle
\noindent Source module: \texttt{Proximal\allowbreak{}Langevin/\allowbreak{}Prox\allowbreak{}Parameter.\allowbreak{}lean}. The accompanying HTML pairs this mathematical exposition with the exact Lean statements and records the scope of the formalization.


For a convex bounded-below \(K\)-Lipschitz function, the actual minimizer satisfies \(\|P_\lambda(x)-P_\eta(x)\|\le K|\lambda-\eta|\). This feeds a pathwise comparison of PIPGLA chains with different positive proximal parameters.


\section{\texorpdfstring{Bounded supports imply value Lipschitzness}{Bounded supports imply value Lipschitzness}}


Choose a genuine subgradient \(s_x\) at each point with \(\|s_x\|\le K\). Its support at \(y\) gives \(f(x)-f(y)\le\langle s_x,x-y\rangle\le K\|x-y\|\). Reverse \(x,y\) to bound the absolute difference. This derives global Lipschitzness from a bounded section, without differentiability.




\textbf{Correspondence with the displayed declarations.}


\begin{itemize}

\item \texttt{lipschitz\_\allowbreak{}of\_\allowbreak{}bounded\_\allowbreak{}subgradients}: A bounded genuine subgradient at every input implies global Lipschitzness. In particular, the paper's uniformly bounded minimal section supplies this hypothesis, without adding a separate Lipschitz assumption to C1.
\end{itemize}


{\raggedright\noindent Lean declarations: \texttt{lipschitz\_\allowbreak{}of\_\allowbreak{}bounded\_\allowbreak{}subgradients}.\par}

\section{\texorpdfstring{Compare two actual proximal minimizers}{Compare two actual proximal minimizers}}


Write \(q=P_\eta(x)\) and \(s=(x-q)/\eta\). Actual proximal optimality gives \(s\in\partial f(q)\) and \(\|s\|\le K\). The same support condition shows that \(q\) is also the \(\lambda\)-proximal point of the shifted input \(q+\lambda s\). Nonexpansiveness at scale \(\lambda\) therefore yields

\[\|P_\lambda(x)-q\|\le\|x-(q+\lambda s)\|=|\eta-\lambda|\|s\|\le K|\lambda-\eta|.\]

This proves the claimed sensitivity for every pair of positive scales.




Let \(q=P_\eta(x)\) and \(s=(x-q)/\eta\). The optimality theorem gives \(s\in\partial f(q)\), and global value Lipschitzness gives \(\|s\|\le K\). Reversing the optimality characterization shows that the same \(q\) is \(P_\lambda(q+\lambda s)\). Nonexpansiveness at parameter \(\lambda\) therefore gives \[\|P_\lambda(x)-P_\eta(x)\|\le\|x-(q+\lambda s)\|=|\eta-\lambda|\|s\|\le K|\lambda-\eta|.\] The local identities \texttt{hq}, \texttt{hxs} and \texttt{hid} implement exactly these substitutions; both parameters must be positive to cancel their residual denominators.




\textbf{Correspondence with the displayed declarations.}


\begin{itemize}

\item \texttt{prox\_\allowbreak{}parameter\_\allowbreak{}sensitivity}: The actual proximal point changes by at most \texttt{K*|lam-\allowbreak{}eta|} when the positive smoothing parameter changes.
\end{itemize}


{\raggedright\noindent Lean declarations: \texttt{prox\_\allowbreak{}parameter\_\allowbreak{}sensitivity}.\par}

\section{\texorpdfstring{The averaged pair map preserves a uniform map difference}{The averaged pair map preserves a uniform map difference}}


Coordinate linearity proves \(\operatorname{Pair}_{P-Q}=\operatorname{Pair}_P-\operatorname{Pair}_Q\). If \(\sup_x\|P(x)-Q(x)\|\le c\), Jensen bounds the output square by the average of quantities at most \(c^2\), so the output norm is at most \(c\). Apply this with the single-pair proximal sensitivity to obtain the same \(K|\lambda-\eta|\) bound for the actual particle map.




Linearity in the supplied pair map gives \(\operatorname{pairwiseMap}(P-Q)=\operatorname{pairwiseMap}(P)-\operatorname{pairwiseMap}(Q)\). If \(\|P(x)-Q(x)\|\le c\) for all pairs, the averaged squared-distance inequality bounds the resulting state difference by \(N^{-1}\sum_i c^2=c^2\). Taking nonnegative square roots proves the same uniform \(c\), with no particle-number loss. Substitute \(c=K|\lambda-\eta|\) from the previous proximal sensitivity estimate.




\textbf{Correspondence with the displayed declarations.}


\begin{itemize}

\item \texttt{pairwise\allowbreak{}Map\_\allowbreak{}sub}: Pairwise averaging is linear in the supplied joint map.
\item \texttt{pairwise\allowbreak{}Map\_\allowbreak{}uniform\_\allowbreak{}difference\_\allowbreak{}le}: A uniform difference of joint maps stays uniformly bounded after the common-parameter averaging and latent-coordinate scaling.
\item \texttt{pairwise\_\allowbreak{}prox\_\allowbreak{}parameter\_\allowbreak{}sensitivity}: The proximal-parameter sensitivity has no particle-number loss.
\end{itemize}


{\raggedright\noindent Lean declarations: \texttt{pairwise\allowbreak{}Map\_\allowbreak{}sub}, \texttt{pairwise\allowbreak{}Map\_\allowbreak{}uniform\_\allowbreak{}difference\_\allowbreak{}le}, \texttt{pairwise\_\allowbreak{}prox\_\allowbreak{}parameter\_\allowbreak{}sensitivity}.\par}

\section{\texorpdfstring{Solve the pathwise parameter-\allowbreak{}perturbation recursion}{Solve the pathwise parameter-perturbation recursion}}


Insert an intermediate update with the same parameter, then combine shared-noise contraction and the uniform parameter difference:

\[e_{n+1}\le(1-m\gamma)e_n+K|\lambda-\eta|.\]

Induction yields \(e_n\le(1-m\gamma)^ne_0+K|\lambda-\eta|\sum_{j<n}(1-m\gamma)^j\). For \(\gamma>0\), the admissible contraction has \(0\le1-m\gamma<1\), so its sum is at most \(1/(m\gamma)\). These bounds hold for every supplied innovation path.




Insert an intermediate update with input \(w\) and proximal parameter \(\lambda\). The triangle inequality separates the state change from the parameter change, giving \[e_{k+1}\le(1-m\gamma)e_k+K|\lambda-\eta|.\] The finite-horizon theorem sums this derived pathwise recurrence, while \(\gamma>0\) gives the strict factor \(1-m\gamma<1\) needed for the floor \(K|\lambda-\eta|/(m\gamma)\). If the two runs start together, the initial term vanishes. No independence hypothesis is needed: both runs use the same arbitrary innovation sequence.




\textbf{Correspondence with the displayed declarations.}


\begin{itemize}

\item \texttt{pipgla\allowbreak{}Update\_\allowbreak{}parameter\_\allowbreak{}comparison}: One step of two actual splitting schemes with different proximal parameters, the same time step, and the same supplied innovation.
\item \texttt{pipgla\allowbreak{}Chain\_\allowbreak{}parameter\_\allowbreak{}comparison\_\allowbreak{}sum}: Exact finite-horizon comparison of two actual splitting recursions.
\item \texttt{pipgla\allowbreak{}Chain\_\allowbreak{}parameter\_\allowbreak{}comparison\_\allowbreak{}bound}: The resulting pathwise difference floor is valid for any two positive proximal parameters; it is not assumed to be an algorithm error bound.
\end{itemize}


{\raggedright\noindent Lean declarations: \texttt{pipgla\allowbreak{}Update\_\allowbreak{}parameter\_\allowbreak{}comparison}, \texttt{pipgla\allowbreak{}Chain\_\allowbreak{}parameter\_\allowbreak{}comparison\_\allowbreak{}sum}, \texttt{pipgla\allowbreak{}Chain\_\allowbreak{}parameter\_\allowbreak{}comparison\_\allowbreak{}bound}.\par}

\section{\texorpdfstring{A proximal corridor gives an explicit scalar bound}{A proximal corridor gives an explicit scalar bound}}


If \(0<m\gamma<1\) and \(\gamma\le\lambda\le\gamma/(1-m\gamma)\), then \(|\lambda-\gamma|=\lambda-\gamma\le m\gamma^2/(1-m\gamma)\). Divide by the positive \(m\gamma\) and multiply by \(K\) to get \(K|\lambda-\gamma|/(m\gamma)\le K\gamma/(1-m\gamma)\). This is a scalar consequence of the true sensitivity theorem.




In the corridor \(\gamma\le\lambda\le\gamma/(1-m\gamma)\), positivity of \(1-m\gamma\) gives \[0\le\lambda-\gamma\le\frac{m\gamma^2}{1-m\gamma}.\] Therefore \(K|\lambda-\gamma|/(m\gamma)\le K\gamma/(1-m\gamma)\). The Lean proof first removes the absolute value using the lower corridor inequality, then clears the two positive denominators. This term measures the additional error incurred by changing the proximal parameter away from the equal-step algorithm.




\textbf{Correspondence with the displayed declarations.}


\begin{itemize}

\item \texttt{prox\_\allowbreak{}parameter\_\allowbreak{}corridor\_\allowbreak{}floor}: In the paper's proximal-parameter corridor, the additional comparison floor against \texttt{eta=gamma} is of order \texttt{gamma}, with every denominator displayed.
\end{itemize}


{\raggedright\noindent Lean declarations: \texttt{prox\_\allowbreak{}parameter\_\allowbreak{}corridor\_\allowbreak{}floor}.\par}

\end{document}
